% LUA101 final exam: all modules, weighted toward modules 8-14; 120 minutes, 100 points.
%
% One source gives both versions. From this folder:
%   latexmk -lualatex final.tex
%     the exam, without the solutions: final.pdf
%   latexmk -lualatex -jobname=final-key -usepretex='\PassOptionsToClass{answers}{exam}' final.tex
%     the key, with every solution and its grading notes: final-key.pdf
% The latexmkrc in this folder adds ../../latex to TEXINPUTS, for the shared preamble. Both PDFs
% stay out of Git. This source holds the key, so it lives in solutions/, never in exams/.
\documentclass[11pt,addpoints]{exam}
\usepackage{preamble}
\usepackage[margin=1in]{geometry}
\usepackage{booktabs}

\pagestyle{headandfoot}
\firstpageheader{LUA101: Lua from Front to Back}{}{Final Exam}
\runningheader{LUA101}{Final Exam}{Page \thepage\ of \numpages}
\firstpagefooter{}{Page \thepage\ of \numpages}{}
\runningfooter{}{}{}

\renewcommand{\solutiontitle}{\noindent\textbf{Solution:}\par\noindent}
\unframedsolutions
\SolutionEmphasis{\color{black!75}}

\lstset{style=lua}
\lstdefinestyle{c}{style=lua, language=C}

% Each question opens with its number and its total points, parts included.
\qformat{\textbf{Question \thequestion}\quad(\totalpoints\ points)\hfill}

\begin{document}

\begin{center}
  {\Large\bfseries LUA101: Lua from Front to Back}\\[0.5ex]
  {\large Final Exam: All Modules}
\end{center}

\vspace{1em}
\noindent Name: \makebox[0.6\textwidth]{\hrulefill}

\vspace{1em}
\noindent\textbf{Instructions.}
\begin{itemize}
  \item You have 120 minutes. The exam has \numquestions\ questions worth \numpoints\ points. Each
    question gives its points, and the learning objectives it assesses; where a question has
    parts, each part gives its share.
  \item Use no computer, phone, or other device. Write each answer in the space provided.
  \item All code is Lua 5.4 unless a question names another version. When a question asks what a
    line prints, write exactly what it prints; where \code|print| separates values with a tab, a
    space is enough.
  \item If an expression or a line raises an error, write ``error'' and say why.
  \item C code uses the Lua 5.4 C API, with \code|lua.h|, \code|lauxlib.h|, and \code|lualib.h|
    included. Small mistakes in C syntax cost nothing if the meaning is clear.
\end{itemize}

\begin{center}
  \gradetable[v][pages]
\end{center}

\newpage
\begin{questions}

% ---------------------------------------------------------------------------------------------
\question \textit{(LO-3.3)} Consider this code.

\begin{lstlisting}
local function shallow(t)
  local c = {}
  for k, v in pairs(t) do c[k] = v end
  return c
end
local a = { n = 1, list = { 1, 2 } }
local b = shallow(a)
b.n = 2
b.list[1] = 10
print(a.n, a.list[1], a.list == b.list)
\end{lstlisting}

\begin{parts}
  \part[3] What does the last line print? One point for each value.
  \begin{solutionorbox}[0.6in]
    \code|1  10  true|. The copy has its own field \code|n|, so changing \code|b.n| leaves
    \code|a.n| alone; but \code|b.list| is the same table as \code|a.list|, because a shallow copy
    copies the reference, so \code|b.list[1] = 10| changes \code|a.list[1]|.
  \end{solutionorbox}
  \part[2] Write \code|deep(t)|, which copies a table and every table inside it, for tables
    without cycles. Other values are returned as they are.
  \begin{solutionorbox}[1.6in]
\begin{lstlisting}
local function deep(t)
  if type(t) ~= "table" then return t end
  local c = {}
  for k, v in pairs(t) do c[k] = deep(v) end
  return c
end
\end{lstlisting}
    One point for recursing on the values, one for returning values that are not tables
    unchanged. Copying the keys with \code|deep| as well, or the metatable, is also correct.
  \end{solutionorbox}
\end{parts}

% ---------------------------------------------------------------------------------------------
\question \textit{(LO-4.4)} Let \code|s| be the five-letter word naïve, written in UTF-8, in
which the letter ï takes two bytes and every other letter one.

\begin{parts}
  \part[1] What is \code|#s|? \hfill \ifprintanswers 6: \code|#| counts bytes. \else \makebox[2in]{\hrulefill} \fi
  \part[1] What is \code|utf8.len(s)|? \hfill \ifprintanswers 5: it counts characters. \else \makebox[2in]{\hrulefill} \fi
  \part[1] What does \code|s:sub(3, 3)| return?
  \begin{solutionorbox}[0.5in]
    A string of one byte, the first of the two bytes of ï, which is not a character on its own:
    \code|string.sub| counts bytes, not characters.
  \end{solutionorbox}
  \part[2] A loop builds a report with \code|out = out .. line .. "\n"| for each of 10{,}000
    lines. Explain why this is slow, and how \code|table.concat| avoids it.
  \begin{solutionorbox}[1in]
    Strings cannot change, so each \code|..| builds a new string and copies everything built so
    far: the total copying grows with the square of the number of lines. Collecting the lines in a
    list and calling \code|table.concat(lines, "\n")| copies each byte once. One point for the
    repeated copying, one for the remedy.
  \end{solutionorbox}
\end{parts}

% ---------------------------------------------------------------------------------------------
\question \textit{(LO-5.1)} Consider this code.

\begin{lstlisting}
local adders, total = {}, 0
for i = 1, 3 do
  adders[i] = function(x) total = total + i * x; return total end
end
local a1 = adders[1](1)
local a3 = adders[3](1)
local a2 = adders[2](10)
print(a1, a3, a2)
\end{lstlisting}

\begin{parts}
  \part[2] What does the last line print?
  \begin{solutionorbox}[0.5in]
    \code|1  4  24|: \code|total| becomes $0 + 1 \cdot 1 = 1$, then $1 + 3 \cdot 1 = 4$, then
    $4 + 2 \cdot 10 = 24$. Two points for all three values; one point for two of them.
  \end{solutionorbox}
  \part[2] Which variables do the three closures share, and which does each have its own copy of?
  \begin{solutionorbox}[0.8in]
    They share \code|total|, one local declared outside the loop. Each has its own \code|i|,
    because each pass of a numeric \code|for| creates a new local \code|i|. One point for each.
  \end{solutionorbox}
\end{parts}

% ---------------------------------------------------------------------------------------------
\question \textit{(LO-7.3)} A program caches the area of each rectangle.

\begin{lstlisting}
local cache = setmetatable({}, { __mode = "k" })
local function area(rect)
  local a = cache[rect]
  if not a then a = rect.w * rect.h; cache[rect] = a end
  return a
end
local r = { w = 2, h = 3 }
print(area(r))
r = nil
collectgarbage()
\end{lstlisting}

\begin{parts}
  \part[1] After the last line, does \code|cache| still hold an entry? Why?
  \begin{solutionorbox}[0.6in]
    No. With \code|__mode = "k"| the keys are weak, and the rectangle is no longer referred to
    from anywhere else, so the full collection removes the entry.
  \end{solutionorbox}
  \part[2] What would change with \code|__mode = "v"| instead, and why?
  \begin{solutionorbox}[0.8in]
    The entry would stay. The keys would then be strong, so the cache itself would keep every
    rectangle alive, and the values are numbers, which the collector never removes from a weak
    table. The cache would grow forever. One point for ``the entry stays'', one for the reason.
  \end{solutionorbox}
  \part[2] Without the call to \code|collectgarbage|, when would the entry disappear?
  \begin{solutionorbox}[0.8in]
    At some collection after the rectangle becomes unreachable: when the collector next finishes
    a cycle that finds it unreachable, which depends on how much the program allocates. Entries
    are not removed at the moment the last strong reference goes. One point for ``at a later
    collection'', one for ``not immediately''.
  \end{solutionorbox}
\end{parts}

% ---------------------------------------------------------------------------------------------
\question \textit{(LO-8.1, LO-8.2)} Consider this code.

\begin{lstlisting}
local Account = {}
Account.__index = Account
function Account.new(owner, balance)
  return setmetatable({ owner = owner, balance = balance or 0 }, Account)
end
function Account:deposit(amount) self.balance = self.balance + amount end
function Account:describe() return self.owner .. ": " .. self.balance end

local Savings = setmetatable({}, { __index = Account })
Savings.__index = Savings
function Savings.new(owner, balance, rate)
  local s = Account.new(owner, balance)
  s.rate = rate
  return setmetatable(s, Savings)
end
function Savings:add_interest() self:deposit(self.balance * self.rate) end
function Savings:describe() return Account.describe(self) .. " (savings)" end

local s = Savings.new("Ada", 100, 0.5)
s:add_interest()
print(s:describe())                                    -- (a)
print(getmetatable(s) == Savings, rawget(s, "deposit")) -- (b)
\end{lstlisting}

\begin{parts}
  \part[1] What does line (a) print?
  \begin{solutionorbox}[0.4in]
    \code|Ada: 150.0 (savings)|: the interest is a float, so the balance becomes \code|150.0|.
  \end{solutionorbox}
  \part[2] What does line (b) print? Explain how \code|s:deposit| is found.
  \begin{solutionorbox}[0.9in]
    \code|true  nil|. \code|deposit| is not stored in \code|s|. Looking it up in \code|s| fails,
    so Lua follows \code|s|'s metatable's \code|__index|, which is \code|Savings|; it is not
    there either, so Lua follows \code|Savings|'s metatable's \code|__index|, which is
    \code|Account|, where it is found. One point for the output, one for the chain.
  \end{solutionorbox}
  \part[1] What happens on \code|Account.new("Bo"):add_interest()|?
  \begin{solutionorbox}[0.5in]
    An error: \code|attempt to call a nil value (method 'add_interest')|. An \code|Account| does
    not inherit from \code|Savings|.
  \end{solutionorbox}
  \part[2] What happens on \code|s.deposit(10)|, and why?
  \begin{solutionorbox}[0.8in]
    An error, \code|attempt to index a number value (local 'self')|. With a dot, 10 is passed as
    \code|self|, and \code|amount| is \code|nil|; the method then tries to read
    \code|self.balance| from a number. The colon in \code|s:deposit(10)| passes \code|s| as
    \code|self|.
  \end{solutionorbox}
  \part[3] A \code|Report| must write its text either to a file or to a list in memory, chosen
    when the report is created, and more destinations may come later. Would you make the
    destinations subclasses of \code|Report|, or give each report a destination object it calls?
    Defend your choice.
  \begin{solutionorbox}[1.4in]
    Composition: give each report a \code|sink| with a \code|write(text)| method, such as a file
    sink and a memory sink. The destination is chosen when the report is created, and new
    destinations need no change to \code|Report|; with subclasses, every kind of report would
    need one subclass per destination, and a report could not change destination. Inheritance
    fits when the subclass is a kind of the parent and adds or changes its behavior, which a
    destination is not. Three points for a clear choice defended with these trade-offs. An answer
    that chooses inheritance earns full marks only if it argues from the same trade-offs, for
    example that there will only ever be two destinations and no other kinds of report.
  \end{solutionorbox}
\end{parts}

% ---------------------------------------------------------------------------------------------
\question \textit{(LO-8.3)} Consider this constructor.

\begin{lstlisting}
local function counter()
  local n = 0
  return {
    inc = function() n = n + 1; return n end,
    get = function() return n end,
  }
end
\end{lstlisting}

\begin{parts}
  \part[1] Why can no code outside \code|counter| change \code|n| except through \code|inc|?
  \begin{solutionorbox}[0.6in]
    \code|n| is a local of \code|counter|; only the two closures created in the call refer to it,
    and it is not a field of the returned table.
  \end{solutionorbox}
  \part[2] A program creates 100{,}000 counters. Compare the memory this design uses with a
    class whose methods are in a metatable shared by every counter.
  \begin{solutionorbox}[1in]
    Each closure-based counter holds its own table, two new closures, and the upvalue
    \code|n|, so the memory per counter is several objects. Metatable-based counters each hold one
    table with a field \code|n|, and all share one copy of the methods. One point for ``two
    closures per counter'', one for ``shared methods''.
  \end{solutionorbox}
  \part[1] Give one thing that is easier with metatable-based counters.
  \begin{solutionorbox}[0.6in]
    Any one of: changing or adding a method for every counter at once, in the class; inheriting
    from the class; or reading and saving a counter's state, which is in its fields.
  \end{solutionorbox}
\end{parts}

% ---------------------------------------------------------------------------------------------
\question[6] \textit{(LO-9.1)} Write everything this code prints, in order.

\begin{lstlisting}
local co = coroutine.create(function(x, y)
  print("start", x, y)
  local c = coroutine.yield(x + y)
  print("got", c)
  local p, q = coroutine.yield(c * 2)
  return p + q
end)
print(coroutine.resume(co, 1, 2))
print(coroutine.status(co))
print(coroutine.resume(co, 10))
print(coroutine.resume(co, 3, 4))
print(coroutine.status(co))
print(coroutine.resume(co))
\end{lstlisting}

\begin{solutionorbox}[2in]
\begin{lstlisting}[language={}]
start  1  2
true  3
suspended
got  10
true  20
true  7
dead
false  cannot resume dead coroutine
\end{lstlisting}
  One point for each of the six \code|print| statements at the end, with what the coroutine
  prints while it runs counted with the \code|resume| that runs it. The first \code|resume| passes
  1 and 2 as arguments; later ones pass their values as the results of \code|yield|; the
  coroutine's return value comes back from the \code|resume| that ends it.
\end{solutionorbox}

% ---------------------------------------------------------------------------------------------
\question \textit{(LO-9.2)} Generators with \code|coroutine.wrap|.

\begin{parts}
  \part[2] Write \code|range(n)|, which returns a generator, for use in a generic \code|for|,
    that yields the integers from 1 to \code|n|.
  \begin{solutionorbox}[1.2in]
\begin{lstlisting}
local function range(n)
  return coroutine.wrap(function()
    for i = 1, n do coroutine.yield(i) end
  end)
end
\end{lstlisting}
  \end{solutionorbox}
  \part[3] Write \code|map(f, gen)|, which returns a new generator that yields \code|f(v)| for
    each value \code|v| that the generator \code|gen| yields.
  \begin{solutionorbox}[1.4in]
\begin{lstlisting}
local function map(f, gen)
  return coroutine.wrap(function()
    for v in gen do coroutine.yield(f(v)) end
  end)
end
\end{lstlisting}
    One point for returning a wrapped coroutine, one for consuming \code|gen| inside it, and one
    for yielding \code|f(v)|.
  \end{solutionorbox}
  \part[1] What does this print?
\begin{lstlisting}
for x in map(function(v) return v * v end, range(3)) do io.write(x, " ") end
\end{lstlisting}
  \begin{solutionorbox}[0.4in]
    \code|1 4 9 |
  \end{solutionorbox}
\end{parts}

% ---------------------------------------------------------------------------------------------
\question \textit{(LO-9.3)} This scheduler runs tasks that sleep. \code|sleep(t)| yields the
number of time units to sleep, and \code|run| always resumes the task that wakes earliest; when two
tasks wake at the same time, it resumes the one that was spawned first.

\begin{lstlisting}
local clock, tasks = 0, {}
local function spawn(fn)
  tasks[#tasks + 1] = { co = coroutine.create(fn), wake = 0 }
end
local function sleep(t) coroutine.yield(t) end
local function run()
  while #tasks > 0 do
    local best = 1
    for i = 2, #tasks do
      if tasks[i].wake < tasks[best].wake then best = i end
    end
    local t = tasks[best]
    clock = t.wake
    local _, delay = coroutine.resume(t.co)
    if coroutine.status(t.co) == "dead" then table.remove(tasks, best)
    else t.wake = clock + delay end
  end
end
spawn(function() for i = 1, 2 do print(clock, "A", i); sleep(3) end end)
spawn(function() for i = 1, 3 do print(clock, "B", i); sleep(2) end end)
run()
\end{lstlisting}

\begin{parts}
  \part[3] Write everything it prints, in order.
  \begin{solutionorbox}[1.4in]
\begin{lstlisting}[language={}]
0  A  1
0  B  1
2  B  2
3  A  2
4  B  3
\end{lstlisting}
    Three points for all five lines in order; two for one line out of place; one for the right
    lines in a wrong order.
  \end{solutionorbox}
  \part[2] Before \code|run()| starts, this line spawns a third task. What happens, and why?
\begin{lstlisting}
spawn(function() while true do end end)
\end{lstlisting}
  \begin{solutionorbox}[0.9in]
    The program hangs once that task runs: it never yields, and a cooperative scheduler can
    switch tasks only when a task yields, so neither of the other tasks runs again. One point for
    ``it hangs and stops the others'', one for the reason.
  \end{solutionorbox}
  \part[1] Name one way to guard a scheduler against such a task.
  \begin{solutionorbox}[0.6in]
    Any one of: run each task under an instruction-count hook that raises an error after too many
    instructions; or review tasks so each loop yields. A pure cooperative scheduler cannot take the
    processor from a task that never yields.
  \end{solutionorbox}
\end{parts}

% ---------------------------------------------------------------------------------------------
\question \textit{(LO-10.2)} A web service counts the lines of a file that a user names:

\begin{lstlisting}
local function count_lines(path)
  return os.execute("wc -l " .. path)
end
\end{lstlisting}

\begin{parts}
  \part[1] Give a value of \code|path| that makes this function delete a file named
    \code|data.db|.
  \begin{solutionorbox}[0.5in]
    \code|x; rm data.db|: the semicolon ends the \code|wc| command and starts another. Any value
    with a shell separator or a substitution works, such as \code|$(rm data.db)|.
  \end{solutionorbox}
  \part[2] Describe two safer ways to count the lines.
  \begin{solutionorbox}[0.9in]
    Count them in Lua, with no shell at all: \code|for _ in io.lines(path) do n = n + 1 end|.
    Or, if a command is needed, quote the path for the shell (in single quotes, with each single
    quote inside written as \code|'\''|), after checking that it names a file the user may read.
    One point each.
  \end{solutionorbox}
  \part[1] What does this expression return?
\begin{lstlisting}
os.date("%Y-%m-%d", os.time({ year = 2026, month = 2, day = 30, hour = 12 }))
\end{lstlisting}
  \begin{solutionorbox}[0.5in]
    \code|"2026-03-02"|: \code|os.time| normalizes February 30 to March 2, because February 2026
    has 28 days.
  \end{solutionorbox}
\end{parts}

% ---------------------------------------------------------------------------------------------
\question \textit{(LO-10.3)} Data that must not run as code.

\begin{parts}
  \part[3] A program saves its settings as a Lua chunk, \code|return { name = "..." }|, and reads
  them back with \code|dofile("settings.lua")|. Explain the risk if another user can edit the
  file, and describe how to read the settings safely.
  \begin{solutionorbox}[1.3in]
    \code|dofile| runs whatever the file contains, with the program's globals, so an edited file
    can call \code|os.execute|, read files, or change the program. A safer reader parses the
    format instead of running it. If the file must stay a Lua chunk, load it in text mode with an
    empty environment, \code|load(text, "=settings", "t", {})|, so that bytecode is refused and
    no library is reachable; but the chunk can still loop forever, so a parser is safest. One point
    for the risk, one for the empty environment and text mode, one for naming what
    \code|load| still allows, or for parsing instead.
  \end{solutionorbox}
  \part[2] Give the fields of this CSV line, in which a quoted field may hold commas and a doubled
    quote stands for one quote: \code|a,"b,c","d""e"|
  \begin{solutionorbox}[0.5in]
    Three fields: \code|a|, \code|b,c|, and \code|d"e|. One point for three fields, one for the
    contents of the second and third.
  \end{solutionorbox}
  \part[1] Why does a program that writes Lua values with \code|string.format("%q", v)| not need
    to escape quotes or newlines itself?
  \begin{solutionorbox}[0.6in]
    \code|%q| writes a string as a Lua literal that reads back as the same string, escaping
    quotes, backslashes, and newlines, and control characters.
  \end{solutionorbox}
\end{parts}

% ---------------------------------------------------------------------------------------------
\question \textit{(LO-12.1, LO-12.3)} Measuring and speeding up.

\begin{parts}
  \part[2] A teammate reports: ``Version B is faster: A took 0.21 s and B took 0.19 s.'' Each
    version ran once, on a laptop with other programs open. Why does the measurement not support
    the claim, and what must the report include?
  \begin{solutionorbox}[1in]
    One run of each cannot separate a difference of 10\% from the noise of a busy machine. Run
    each version several times, collecting garbage before each run, and report the minimum or
    the median, with the input, the hardware, the operating system, and the Lua version. One
    point for the reason, one for what to report.
  \end{solutionorbox}
  \part[2] This function is a program's hot path, called with thousands of lines and a stop list
    of 200 words. Rewrite it to be faster, and say why the change helps.
\begin{lstlisting}
local function count_words(lines, stop)
  local n = 0
  for _, line in ipairs(lines) do
    for word in line:gmatch("%a+") do
      local skip = false
      for _, s in ipairs(stop) do
        if s == word then skip = true end
      end
      if not skip then n = n + 1 end
    end
  end
  return n
end
\end{lstlisting}
  \begin{solutionorbox}[1.6in]
    Build a set of the stop words once, and look each word up in it:
\begin{lstlisting}
local is_stop = {}
for _, s in ipairs(stop) do is_stop[s] = true end
-- in the loop:
if not is_stop[word] then n = n + 1 end
\end{lstlisting}
    Each word then costs one table lookup instead of 200 comparisons. One point for the set, one
    for the reason.
  \end{solutionorbox}
  \part[1] How would you show that the change was worth making?
  \begin{solutionorbox}[0.6in]
    Time both versions on the same, realistic input, several times each, check that they return
    the same count, and report both times with their conditions.
  \end{solutionorbox}
\end{parts}

% ---------------------------------------------------------------------------------------------
\question \textit{(LO-12.2)} How Lua stores tables and strings.

\begin{parts}
  \part[2] Compare the cost of \code|t[#t + 1] = x| with the cost of
    \code|table.insert(t, 1, x)|, for a list of 100{,}000 items.
  \begin{solutionorbox}[0.8in]
    Appending costs about the same however long the list is (the array part grows by doubling, so
    its occasional resizes average out). Inserting at position 1 moves every item up one place, so
    its cost grows with the length of the list: here, 100{,}000 moves for each insertion.
  \end{solutionorbox}
  \part[1] In \code|local t = { 10, 20, x = 1 }|, where are the values stored?
  \begin{solutionorbox}[0.5in]
    10 and 20 in the array part, at indices 1 and 2; \code|x| in the hash part.
  \end{solutionorbox}
  \part[2] Why can comparing two equal strings of 1{,}000 characters with \code|==| take longer
    than comparing two equal strings of 10 characters?
  \begin{solutionorbox}[0.9in]
    Lua 5.4 interns short strings, of 40 bytes or fewer: it keeps one copy of each, so \code|==|
    compares two pointers. Long strings are not interned, so two equal long strings made
    separately are different objects, and \code|==| compares them byte by byte. One point for
    interning, one for the byte-by-byte comparison.
  \end{solutionorbox}
\end{parts}

% ---------------------------------------------------------------------------------------------
\question[6] \textit{(LO-13.1)} The stack holds \code|'a' 'b'|, bottom first. Give the stack after
each call, bottom first. Each call starts from the stack the previous one left.

\begin{center}
\renewcommand{\arraystretch}{1.4}
\begin{tabular}{@{}l l p{2.6in}@{}}
  \toprule
  & Call & Stack afterward \\
  \midrule
  (a) & \code|lua_pushinteger(L, 1)| & \ifprintanswers \code|'a' 'b' 1| \fi \\
  (b) & \code|lua_pushvalue(L, -2)| & \ifprintanswers \code|'a' 'b' 1 'b'| \fi \\
  (c) & \code|lua_insert(L, 1)| & \ifprintanswers \code|'b' 'a' 'b' 1| \fi \\
  (d) & \code|lua_remove(L, 2)| & \ifprintanswers \code|'b' 'b' 1| \fi \\
  (e) & \code|lua_settop(L, 4)| & \ifprintanswers \code|'b' 'b' 1 nil| \fi \\
  (f) & \code|lua_replace(L, -3)| & \ifprintanswers \code|'b' nil 1| \fi \\
  \bottomrule
\end{tabular}
\end{center}

\begin{solution}
One point for each row (6 points). \code|pushvalue(L, -2)| copies the value second from the top,
\code|'b'|. \code|insert(L, 1)| moves the top value to the bottom. \code|settop| fills new places
with \code|nil|. \code|replace(L, -3)| pops the top value, \code|nil|, into the place that was
third from the top before the pop, index 2. A C program that prints the stack after each call
confirms every row.
\end{solution}

% ---------------------------------------------------------------------------------------------
\question[6] \textit{(LO-13.2)} Write the C function \code|l_count| of a module \code|strx|, so
that \code|strx.count(s, c)| returns the number of times the one-character string \code|c|
occurs in the string \code|s|. A wrong argument must raise an error that names it; in
particular, \code|strx.count("banana", "an")| must raise this error:
\begin{lstlisting}[language={}]
bad argument #2 to 'count' (one character expected)
\end{lstlisting}
Also write the list of functions and the \code|luaopen_strx| function that
\code|require("strx")| calls.

\begin{solutionorbox}[3.6in]
\begin{lstlisting}[style=c]
static int l_count(lua_State *L) {
  size_t len, clen;
  const char *s = luaL_checklstring(L, 1, &len);
  const char *c = luaL_checklstring(L, 2, &clen);
  luaL_argcheck(L, clen == 1, 2, "one character expected");
  lua_Integer n = 0;
  for (size_t i = 0; i < len; i++) {
    if (s[i] == c[0]) n++;
  }
  lua_pushinteger(L, n);
  return 1;
}

static const luaL_Reg functions[] = {
  { "count", l_count },
  { NULL, NULL },
};

int luaopen_strx(lua_State *L) {
  luaL_newlib(L, functions);
  return 1;
}
\end{lstlisting}
  Two points for checking both arguments, with \code|luaL_checklstring|, and for checking the
  length of \code|c|, with \code|luaL_argcheck|; two for the count, pushed with \code|lua_pushinteger| and with
  \code|return 1|; two for the list, ended by \code|{ NULL, NULL }|, and \code|luaopen_strx| with
  \code|luaL_newlib|. Using the length from \code|luaL_checklstring|, rather than
  \code|strlen|, also counts characters after a zero byte; either earns full marks. Compiled and
  run, this code returns 3 for \code|strx.count("banana", "a")| and raises the required error.
\end{solutionorbox}

% ---------------------------------------------------------------------------------------------
\question \textit{(LO-13.3)} A C program passes each event to a script's function
\code|handle|, and prints what it returns.

\begin{lstlisting}[style=c]
static void handle_all(lua_State *L, const char **events, int n) {
  for (int i = 0; i < n; i++) {
    lua_getglobal(L, "handle");
    lua_pushstring(L, events[i]);
    lua_call(L, 1, 1);
    printf("%s\n", lua_tostring(L, -1));
  }
}
\end{lstlisting}

\begin{parts}
  \part[3] Name three problems in this loop, and what each causes.
  \begin{solutionorbox}[1.4in]
    (1) \code|lua_call| is unprotected: an error in \code|handle| reaches the panic function,
    which aborts the whole program. (2) The result is never popped, so the stack grows by one
    value per event, until it overflows after about a million. (3) \code|lua_tostring| returns
    \code|NULL| when the result is neither a string nor a number, and passing \code|NULL| to
    \code|printf| for \code|%s| is undefined behavior. One point each.
  \end{solutionorbox}
  \part[2] Rewrite the body of the loop to fix all three.
  \begin{solutionorbox}[1.8in]
\begin{lstlisting}[style=c]
lua_getglobal(L, "handle");
lua_pushstring(L, events[i]);
if (lua_pcall(L, 1, 1, 0) != LUA_OK) {
  fprintf(stderr, "handle: %s\n", lua_tostring(L, -1));
} else {
  const char *text = lua_tostring(L, -1);
  printf("%s\n", text ? text : "(no text)");
}
lua_pop(L, 1);   /* the result, or the error message */
\end{lstlisting}
    One point for \code|lua_pcall| with its status checked, one for the single \code|lua_pop| on
    both paths with the \code|NULL| case handled.
  \end{solutionorbox}
\end{parts}

% ---------------------------------------------------------------------------------------------
\question \textit{(LO-13.4)} A module exposes a C struct \code|Image|, whose field \code|pixels|
points to memory from \code|malloc|, as full userdata with a metatable registered as
\code|"Image"|.

\begin{parts}
  \part[1] Why does its \code|__gc| set \code|img->pixels = NULL| after freeing it?
  \begin{solutionorbox}[0.6in]
    So that a second call cannot free the memory twice: a program can call \code|__gc| or a
    \code|close| method itself, and the collector calls \code|__gc| again later.
  \end{solutionorbox}
  \part[1] Why do its methods get the struct with \code|luaL_checkudata(L, 1, "Image")|, and not
    with \code|lua_touserdata(L, 1)|?
  \begin{solutionorbox}[0.6in]
    \code|luaL_checkudata| checks that the value is a userdata with the \code|Image| metatable,
    and raises an error naming the argument otherwise; \code|lua_touserdata| accepts any
    userdata, such as a file handle, and the method would misread its memory.
  \end{solutionorbox}
  \part[1] A loop creates many images and keeps none, yet the process grows to gigabytes before
    the collector frees any. Why?
  \begin{solutionorbox}[0.6in]
    The collector paces itself by the memory Lua allocates, and sees only each userdata's small
    struct, not the pixels that \code|malloc| allocated, so it runs too rarely.
  \end{solutionorbox}
\end{parts}

% ---------------------------------------------------------------------------------------------
\question \textit{(LO-14.1)} Consider this function, written for Lua 5.4.

\begin{lstlisting}
local function halves(list)
  local out = {}
  for i = 1, #list do out[i] = list[i] // 2 end
  return table.unpack(out)
end
\end{lstlisting}

\begin{parts}
  \part[2] Give two reasons it does not run on Lua 5.1.
  \begin{solutionorbox}[0.8in]
    \code|//| does not exist in Lua 5.1, so the file fails to compile, before any line runs; and
    \code|table.unpack| is \code|nil| in 5.1, where the function is the global \code|unpack|. One
    point each.
  \end{solutionorbox}
  \part[3] Rewrite it so that one file runs on Lua 5.1, LuaJIT, Lua 5.4, and Lua 5.5, and gives
    5.4's results on 5.4. Detect features; do not test \code|_VERSION|.
  \begin{solutionorbox}[2in]
\begin{lstlisting}
local unpack = table.unpack or unpack
local chunk = (loadstring or load)("return function(a, b) return a // b end")
local idiv = chunk and chunk() or function(a, b) return math.floor(a / b) end

local function halves(list)
  local out = {}
  for i = 1, #list do out[i] = idiv(list[i], 2) end
  return unpack(out)
end
\end{lstlisting}
    One point for \code|table.unpack or unpack|, one for keeping \code|//| out of the file's own
    text (compiling it from a string), and one for a fallback that rounds down. On each of the
    four Luas, \code|halves({ 7, -7, 4 })| returns \code|3 -4 2|.
  \end{solutionorbox}
\end{parts}

% ---------------------------------------------------------------------------------------------
\question \textit{(LO-14.2)} A program runs plugins in a sandbox: an environment built from an
allowlist, \code|load| in text mode, and an instruction-count hook in the plugin's own coroutine.

\begin{parts}
  \part[1] The allowlist includes \code|load|. Write one line of plugin code that reaches the
    host's real \code|os| library.
  \begin{solutionorbox}[0.4in]
    \code|local os = load("return os")()|: a chunk loaded without an environment gets the real
    global table.
  \end{solutionorbox}
  \part[1] When the instruction limit is reached, the hook raises an error, and from then on
    raises it again at every instruction. Why is raising it once not enough?
  \begin{solutionorbox}[0.6in]
    The plugin can catch one error with \code|pcall| and keep going, as
    \code|while true do pcall(f) end| does. Raising at every instruction makes the code outside
    the \code|pcall| fail too.
  \end{solutionorbox}
  \part[2] Name two kinds of work that the hook cannot stop, and for each, a remedy outside Lua.
  \begin{solutionorbox}[1.2in]
    Any two, one point each with its remedy: one huge allocation in one C call, such as
    \code|("x"):rep(1e9)| (an allocator that refuses to allocate past a limit); a long
    computation in one C call, such as a pattern that backtracks (running the plugin in a separate
    process that is killed after a deadline); compiled LuaJIT code, which does not call the count
    hook (turning the compiler off with \code|jit.off()|, or a separate process).
  \end{solutionorbox}
  \part[1] Why does the sandbox load plugins in text mode?
  \begin{solutionorbox}[0.5in]
    To refuse precompiled bytecode: Lua does not check that a binary chunk's instructions make
    sense, and crafted bytecode can crash the interpreter.
  \end{solutionorbox}
\end{parts}

\end{questions}

\end{document}
